\scanpage{357}
\mspage{96}

Comme \deleted{\unsure{toute}}
\[
\varphi:\tau\mathrel{\unsure{-|\to}}\tau'
\]
induit
\[
\varphi:I_\tau\mathrel{-}I_{\tau'}
\]
transformant flèches rétrogrades en flèches identiques ou \unsure{régtrogrades|rétrogrades}, i.e.
\[
\varphi(\Sigma_\tau)\subset\Sigma_{\tau'},
\]
on trouve le premier des deux foncteurs $(*)$ annoncés p. 94, qui est associé par transposition au foncteur $\gamma$
\diagram{\begin{tikzcd}[ampersand replacement=\&,column sep=small,row sep=large]
(*)\&\underline{\mathrm{Hom}}(\tau,\tau')\times
\underbrace{\underline{\mathrm{Ch}}_{\Sigma,\tau'}(X)}_{\simeq\underline{\mathrm{Hom}}_{\Sigma_{\tau'},\Sigma}(I_{\tau'},X)}
\arrow[d,"\mathrm{can}"']\arrow[r,"\gamma"]\&
\underline{\mathrm{Ch}}_{\Sigma,\tau}(X)\arrow[d,phantom,"\parallel"]\\
\&\underline{\mathrm{Hom}}_{\Sigma_\tau,\Sigma_{\tau'}}(I_\tau,I_{\tau'})\times\underline{\mathrm{Hom}}_{\Sigma_{\tau'},\Sigma}(I_{\tau'},X)
\arrow[r,dotted,"\gamma_0"]\&\underline{\mathrm{Hom}}_{\Sigma_\tau,\Sigma}(I_\tau,X)
\end{tikzcd}}
\unsure{déduit} du foncteur de composition $\gamma_0$ dans ce diagramme : $\gamma=\gamma_0\circ\mathrm{can}$. D'autre part, ce foncteur de composition $\gamma$ induit visiblement un foncteur sur les sous-catégories de chemins à \deleted{\unsure{origine}} \inserted{source} et \deleted{\unsure{extrémité}} \inserted{but} fixés
\[
\underline{\mathrm{Hom}}(\tau,\tau')\times\underline{\mathrm{Ch}}_{\Sigma,\tau'}(X;a,b)\mathrel{-}\underline{\mathrm{Ch}}_{\Sigma,\tau}(X;a,b),
\]
d'où la deuxième des deux flèches annoncées dans $(*)$ p. 94.

Cela montre donc que si
\diagram{\begin{tikzcd}[ampersand replacement=\&,column sep=small]
\varphi_1,\varphi_2:\&\tau\arrow[r,shift left=.5ex,no head]\arrow[r,shift right=.5ex,no head]\&\tau'
\end{tikzcd}}
sont deux morphismes de $\tau$ dans $\tau'$ homotopes dans $\underline{\mathrm{Hom}}(\tau,\tau')$, alors les \deleted{deux} foncteurs induits entre \deleted{\unsure{chemins}} \inserted{cat. de} $\Sigma$-chemins ou \inserted{entre cat. de} $\Sigma$-chemins à \unsure{extrémités} \unsure{précisées|précisés} :

% archive sheet number: 356
